%======================================================================
%  preamble.tex —— 讲义公共样式（版式、字体、配色、盒子、题目排版）
%  用 XeLaTeX 编译；教师版 / 学生版由 \StudentVersion 是否定义来切换
%======================================================================

%---------------------------- 版本开关 ----------------------------
\newif\ifteacher
\ifdefined\StudentVersion \teacherfalse \else \teachertrue \fi

%---------------------------- 版面 ----------------------------
\usepackage[a4paper,left=2.1cm,right=2.1cm,top=2.5cm,bottom=2.3cm,
  headheight=16pt,headsep=0.75cm,footskip=0.95cm]{geometry}

%---------------------------- 字体 ----------------------------
\usepackage[math-style=ISO]{unicode-math}
\setmainfont{TeX Gyre Termes}
\setsansfont{TeX Gyre Heros}[Scale=0.92]
\setmathfont{TeX Gyre Termes Math}
\setmathfont{TeX Gyre Termes Math}[version=bold,FakeBold=1.6]
\setCJKmainfont{Noto Serif CJK SC}[ItalicFont={Noto Serif CJK SC}]
\setCJKsansfont{Noto Sans CJK SC}[BoldFont={Noto Sans CJK SC Bold},ItalicFont={Noto Sans CJK SC}]
\setCJKmonofont{Noto Sans CJK SC}
\setCJKfamilyfont{zhkai}{FandolKai-Regular.otf}
\setCJKfamilyfont{zhfs}{FandolFang-Regular.otf}
\setCJKfamilyfont{zhhei}{Noto Sans CJK SC}
\providecommand{\kaishu}{}\renewcommand{\kaishu}{\CJKfamily{zhkai}}
\providecommand{\fangsong}{}\renewcommand{\fangsong}{\CJKfamily{zhfs}}
\providecommand{\heiti}{}\renewcommand{\heiti}{\sffamily}
% 罗马数字、带圈数字、★ 等符号交给中文字体
\xeCJKDeclareCharClass{CJK}{"2160 -> "217F, "2460 -> "2473, "2605, "2606, "25A0 -> "25FF}
% 常数 π 用正体；“变化量”Δ 用正体
\AtBeginDocument{\renewcommand{\pi}{\symup{π}}}
\newcommand{\dd}{\symup{\Delta}}

%---------------------------- 常用宏包 ----------------------------
\usepackage{amsmath}
\usepackage{xcolor}
\usepackage{graphicx}
\usepackage{array,booktabs,tabularx,makecell,multirow}
\usepackage{enumitem}
\usepackage{multicol}
\usepackage{environ}
\usepackage{caption}
\usepackage{pifont}
\usepackage{fancyhdr}
\usepackage{siunitx}
\sisetup{
  inter-unit-product = \ensuremath{\cdot},
  per-mode = symbol,
  mode = match,
  output-decimal-marker = {.},
}
\usepackage{tikz}
\usetikzlibrary{arrows.meta,calc,decorations.markings,decorations.pathmorphing,
  decorations.pathreplacing,patterns,positioning,shapes.geometric,shadings,
  fadings,intersections,backgrounds,fit}
\usepackage[most]{tcolorbox}
\usepackage[hidelinks,bookmarksnumbered]{hyperref}

%---------------------------- 配色 ----------------------------
\definecolor{mainblue}{HTML}{1F4E79}
\definecolor{lightblue}{HTML}{EEF4FB}
\definecolor{exred}{HTML}{A93226}
\definecolor{warnorange}{HTML}{C0561A}
\definecolor{lightorange}{HTML}{FDF3EA}
\definecolor{extgreen}{HTML}{1E7B4F}
\definecolor{lightgreen}{HTML}{ECF7F0}
\definecolor{histbrown}{HTML}{7A5530}
\definecolor{lightbrown}{HTML}{F8F2EA}
\definecolor{methodpurple}{HTML}{5B2C83}
\definecolor{lightpurple}{HTML}{F4EEF9}
\definecolor{expteal}{HTML}{116A75}
\definecolor{lightteal}{HTML}{EAF5F6}
\definecolor{ansred}{HTML}{C0392B}
\definecolor{gridgray}{HTML}{8A8A8A}
% 物理作图配色
\colorlet{Bline}{blue!65!black}      % 磁感线、磁场符号
\colorlet{Eline}{red!72!black}       % 感生电场线
\colorlet{Icur}{orange!92!red}       % 感应电流、涡流
\colorlet{Fforce}{green!45!black}    % 力
\colorlet{Vvel}{violet!80!black}     % 速度
\colorlet{metal}{gray!28}            % 一般金属
\colorlet{alu}{gray!14}              % 铝
\colorlet{copper}{orange!45!brown!70}% 铜
\colorlet{ironc}{gray!45}            % 铁芯

%---------------------------- 页眉页脚 ----------------------------
\pagestyle{fancy}
\fancyhf{}
\fancyhead[L]{\small\color{mainblue}\sffamily 高中物理 · 电磁感应专题}
\fancyhead[R]{\small\color{mainblue}\sffamily 涡流、电磁阻尼和电磁驱动
  \ifteacher（教师版）\else（学生版）\fi}
\fancyfoot[C]{\small\sffamily — \thepage\ —}
\renewcommand{\headrule}{{\color{mainblue}\hrule height 0.6pt}}
\fancypagestyle{plain}{\fancyhf{}\fancyfoot[C]{\small\sffamily — \thepage\ —}%
  \renewcommand{\headrule}{}}

%---------------------------- 标题格式 ----------------------------
\ctexset{
  section = {
    format      = \Large\sffamily\bfseries\color{mainblue}\raggedright,
    name        = {,、},
    number      = \chinese{section},
    aftername   = {},
    beforeskip  = 2.2ex plus .4ex minus .2ex,
    afterskip   = 1.2ex plus .2ex,
    aftertitle  = {\par\nobreak\vspace{-0.9ex}%
      {\color{mainblue}\rule{\linewidth}{1.1pt}}\par\nobreak\vspace{-0.6ex}},
  },
  subsection = {
    format      = \large\sffamily\bfseries\color{mainblue!88!black}\raggedright,
    name        = {,．},
    number      = \arabic{subsection},
    aftername   = {},
    beforeskip  = 1.6ex plus .3ex minus .2ex,
    afterskip   = 0.8ex plus .1ex,
  },
  subsubsection = {
    format      = \normalsize\sffamily\bfseries\raggedright,
    number      = （\arabic{subsubsection}）,
    aftername   = {},
    beforeskip  = 1.0ex plus .2ex,
    afterskip   = 0.5ex,
  },
}
\setcounter{secnumdepth}{3}
\setcounter{tocdepth}{1}
\captionsetup{font=small,labelsep=quad,skip=4pt,hypcap=false}

%---------------------------- 列表 ----------------------------
\setlist{nosep,leftmargin=2em}
\setlist[enumerate,1]{label=\arabic*．}
\newlist{steps}{enumerate}{1}
\setlist[steps]{label=\ding{\numexpr171+\value{stepsi}\relax},leftmargin=1.8em,nosep}

%---------------------------- 行内标签 ----------------------------
\newcommand{\kd}[1]{{\sffamily\bfseries\color{white}\colorbox{mainblue}{\,#1\,}}}
\newcommand{\biao}[2][mainblue]{{\sffamily\small\bfseries\color{#1}【#2】}}
\newcommand{\key}[1]{{\sffamily\bfseries\color{mainblue}#1}}
\newcommand{\warn}[1]{{\bfseries\color{warnorange}#1}}
\newcommand{\xing}{{\color{warnorange}★}}
\newcommand{\note}[1]{{\kaishu #1}}

%---------------------------- 盒子 ----------------------------
\tcbset{
  handoutbase/.style={enhanced,breakable,arc=2pt,boxrule=0.6pt,
    left=7pt,right=7pt,top=6pt,bottom=6pt,
    before skip=8pt plus 2pt,after skip=8pt plus 2pt,
    fonttitle=\sffamily\bfseries\small},
  sidebar/.style 2 args={handoutbase,colback=#2,colframe=#1,
    boxrule=0pt,leftrule=3pt,arc=0pt,outer arc=0pt,
    coltitle=#1,colbacktitle=#2,
    title style={},
    titlerule=0pt,toptitle=3pt,bottomtitle=0pt},
}
% 知识梳理
\newtcolorbox{zhishi}[1][核心知识]{sidebar={mainblue}{lightblue},title={\ding{43}\ #1}}
% 易错警示
\newtcolorbox{yicuo}[1][易错警示]{sidebar={warnorange}{lightorange},title={\ding{115}\ #1}}
% 拓展提升
\newtcolorbox{tuozhan}[1][拓展提升]{sidebar={extgreen}{lightgreen},title={★\ #1}}
% 物理学史
\newtcolorbox{wulishi}[1][物理学史]{sidebar={histbrown}{lightbrown},title={\ding{46}\ #1}}
% 方法提炼
\newtcolorbox{fangfa}[1][方法提炼]{sidebar={methodpurple}{lightpurple},title={\ding{51}\ #1}}
% 实验探究
\newtcolorbox{shiyan}[1][实验观察]{sidebar={expteal}{lightteal},title={\ding{96}\ #1}}

% 例题
\NewTColorBox[auto counter]{liti}{ o }{handoutbase,
  colback=white,colframe=exred!75!black,boxrule=0.6pt,
  top=10pt,
  title={例\,\thetcbcounter\IfValueT{#1}{\hspace{0.7em}{\mdseries\small（#1）}}},
  attach boxed title to top left={xshift=8pt,yshift=-\tcboxedtitleheight/2},
  boxed title style={colback=exred!75!black,colframe=exred!75!black,
    boxrule=0pt,arc=1.5pt,left=5pt,right=5pt,top=1.5pt,bottom=1.5pt},
}

% 解析（教师版显示；学生版留出空白）
\newcommand{\jxlabel}[1]{{\sffamily\bfseries\small\color{exred}【#1】}}
\newcommand{\jxrule}{\par\vspace{3pt}\noindent
  \tikz\draw[exred!45,dash pattern=on 3pt off 2pt,line width=0.4pt](0,0)--(\linewidth,0);%
  \par\vspace{1pt}}
% 学生版留白用不可丢弃的支撑盒：若用纯 \vspace，盒子恰在留白处分页时
% 剩余部分会被丢弃，连同其中的颜色出栈指令，导致后文颜色错乱
\NewEnviron{jiexi}[1][2.5cm]{%
  \ifteacher
    \jxrule\noindent\jxlabel{解析}\BODY\par
  \else
    \par\noindent\rule{0pt}{#1}\par
  \fi}
\newcommand{\daan}[1]{\ifteacher\par\noindent\jxlabel{答案}#1\par\fi}
% 选择题题干末尾的括号：教师版填答案
\newcommand{\kuohao}[1]{\ifteacher（\,{\color{ansred}\bfseries #1}\,）\else（\qquad）\fi}
% 填空题：\tk[宽度]{答案}，教师版显示答案，学生版留空
\newcommand{\tk}[2][4em]{\ifteacher\underline{\makebox[#1]{\color{ansred}\bfseries #2}}%
  \else\underline{\makebox[#1]{}}\fi}
% 学生版计算题后的作答留白
\newcommand{\liukong}[1]{\ifteacher\else\par\vspace{#1}\fi}
% 判断题：\pd{\dui 或 \cuo}{理由}，教师版给出判断和理由，学生版留空
\newcommand{\dui}{\ding{51}}
\newcommand{\cuo}{\ding{55}}
\newcommand{\pd}[2]{\ifteacher\hfill（\,{\color{ansred}\bfseries #1}\,）\par
  {\small\kaishu\color{black!78}#2}\else\hfill（\qquad）\fi}

%---------------------------- 选择题选项自动排版 ----------------------------
\newlength\xxlen
\newlength\xxmax
\newcommand{\xxmeasure}[1]{\settowidth\xxlen{#1}\ifdim\xxlen>\xxmax\setlength\xxmax{\xxlen}\fi}
\newcommand{\xxitem}[2]{\makebox[#1][l]{#2}}
\newcommand{\xx}[4]{\par\nobreak\vspace{1pt}%
  \setlength\xxmax{0pt}%
  \xxmeasure{A．#1}\xxmeasure{B．#2}\xxmeasure{C．#3}\xxmeasure{D．#4}%
  \ifdim\xxmax<0.23\linewidth
    \noindent\xxitem{0.25\linewidth}{A．#1}\xxitem{0.25\linewidth}{B．#2}%
      \xxitem{0.25\linewidth}{C．#3}\xxitem{0.25\linewidth}{D．#4}\par
  \else\ifdim\xxmax<0.48\linewidth
    \noindent\xxitem{0.5\linewidth}{A．#1}\xxitem{0.5\linewidth}{B．#2}\par
    \noindent\xxitem{0.5\linewidth}{C．#3}\xxitem{0.5\linewidth}{D．#4}\par
  \else
    {\setlength\parindent{0pt}\hangindent=1.6em\hangafter=1 A．#1\par
     \hangindent=1.6em\hangafter=1 B．#2\par
     \hangindent=1.6em\hangafter=1 C．#3\par
     \hangindent=1.6em\hangafter=1 D．#4\par}%
  \fi\fi}

%---------------------------- 作图通用样式 ----------------------------
\tikzset{
  >={Stealth[length=5pt,width=3.6pt]},
  every picture/.style={line cap=round,line join=round,font=\small},
  % 磁场符号：垂直纸面向外 / 向里
  Bout/.pic={\draw[Bline,line width=0.5pt] (0,0) circle (2.3pt); \fill[Bline] (0,0) circle (0.8pt);},
  Bin/.pic={\draw[Bline,line width=0.5pt] (-2.1pt,-2.1pt)--(2.1pt,2.1pt) (-2.1pt,2.1pt)--(2.1pt,-2.1pt);},
  % 电流符号（导线截面）：流出 / 流入纸面
  Iout/.pic={\draw[Icur,line width=0.5pt,fill=white] (0,0) circle (3pt); \fill[Icur] (0,0) circle (1pt);},
  Iin/.pic={\draw[Icur,line width=0.5pt,fill=white] (0,0) circle (3pt);
    \draw[Icur,line width=0.5pt] (-2pt,-2pt)--(2pt,2pt) (-2pt,2pt)--(2pt,-2pt);},
  % 沿路径中点加箭头
  midarrow/.style={postaction={decorate},decoration={markings,
    mark=at position #1 with {\arrow{>}}}},
  midarrow/.default=0.55,
  revarrow/.style={postaction={decorate},decoration={markings,
    mark=at position #1 with {\arrowreversed{>}}}},
  revarrow/.default=0.55,
  % 箭头样式
  vel/.style={->,Vvel,thick},
  force/.style={->,Fforce,thick},
  cur/.style={->,Icur,thick},
  bline/.style={Bline,line width=0.6pt},
  eline/.style={Eline,line width=0.6pt},
  ground/.style={pattern=north east lines,pattern color=gray},
  lbl/.style={font=\footnotesize},
}
% 双极脉冲（作电流-时间图用）：pulse(t, t0, 幅度, 宽度)
\pgfmathdeclarefunction{pulse}{4}{\pgfmathparse{-2*#3*((#1-#2)/#4)*exp(-((#1-#2)/#4)^2)}}
% 在矩形区域内画磁场符号阵列：\Bfield{out|in}{x0}{y0}{x1}{y1}{dx}{dy}
\newcommand{\Bfield}[7]{%
  \pgfmathtruncatemacro\bnx{round((#4-(#2))/(#6))}%
  \pgfmathtruncatemacro\bny{round((#5-(#3))/(#7))}%
  \foreach \bi in {0,...,\bnx}{\foreach \bj in {0,...,\bny}{%
    \pic at ({#2+\bi*(#6)},{#3+\bj*(#7)}) {B#1};}}}

%---------------------------- 其他 ----------------------------
\setlength{\parindent}{2em}
\linespread{1.06}
\renewcommand{\arraystretch}{1.25}
\newcolumntype{Y}{>{\raggedright\arraybackslash}X}
\newcolumntype{C}{>{\centering\arraybackslash}X}
